\section{Sinh lân cận bằng phép nhiễu}
\label{sec:ch03-sinh-lan-can}

\index{phép nhiễu}%
Sau khi đã chọn $x'$, LIME phải tạo dữ liệu để học $g$. Nó sinh các
\tn{phép nhiễu}{perturbation} $z'$ quanh $x'$, dựng lại đầu vào gốc tương ứng
$z$, rồi lấy $f(z)$ làm nhãn để khớp mô hình thay thế. Trong mô tả của bài báo,
với một $x'$ nhị phân, LIME lấy ngẫu nhiên các phần tử đang có mặt để tạo
$z'$; số phần tử được lấy cũng thay đổi giữa các lần sinh
mẫu~\autocite{p09lime}.

Hình~\ref{fig:ch03-quy-trinh-lime} tách những bước này ra. Mũi tên từ $z'$ sang
$z$ là nơi một bit vắng mặt trở thành một can thiệp cụ thể: bỏ từ trong văn
bản, làm mờ vùng ảnh, hoặc một cách tái tạo khác. Mũi tên từ $z$ sang $f(z)$
không mở mô hình \tn{hộp đen}{black box}; nó chỉ hỏi mô hình đã phản ứng thế
nào với đầu vào
đã dựng lại.

\begin{figure}[htbp]
  \centering
  \begin{tikzpicture}[
  box/.style={draw=black!65, rounded corners=2pt, align=center,
    minimum width=4.1cm, minimum height=0.95cm, fill=black!5},
  data/.style={draw=black!65, rounded corners=2pt, align=center,
    minimum width=4.1cm, minimum height=0.85cm, fill=black!15},
  arrow/.style={-{Stealth[length=2mm]}, draw=black!65, thick}
]
  \node[data] (input) {Đầu vào $x$ và biểu diễn $x'$};
  \node[box, below=8mm of input] (sample) {Sinh các mẫu nhiễu $z'$};
  \node[box, below=8mm of sample] (restore) {Dựng lại đầu vào $z$};
  \node[box, below=8mm of restore] (blackbox) {Mô hình hộp đen trả $f(z)$};
  \node[data, right=14mm of blackbox] (weight) {Trọng số\\ lân cận $\pi_x(z)$};
  \node[box, below=10mm of blackbox, xshift=29mm] (surrogate)
    {Khớp mô hình thay thế thưa $g$};

  \draw[arrow] (input) -- (sample);
  \draw[arrow] (sample) -- (restore);
  \draw[arrow] (restore) -- (blackbox);
  \draw[arrow] (blackbox) -- (surrogate);
  \draw[arrow] (weight.south) |- (surrogate.east);
\end{tikzpicture}

  \caption{Quy trình LIME quanh đầu vào $x$. Biểu diễn $x'$ quyết định những
  thành phần có thể bật hoặc tắt; các mẫu $z'$ được dựng lại thành đầu vào $z$
  trước khi mô hình hộp đen trả $f(z)$. Trọng số lân cận và phép khớp tạo ra
  mô hình thay thế $g$, chứ không truy xuất cơ chế bên trong của
  $f$~\autocite{p09lime}.}
  \label{fig:ch03-quy-trinh-lime}
\end{figure}

Không phải mọi $z$ tạo được đều là một đầu vào có nghĩa. Bỏ nhiều từ có thể
biến một câu thành chuỗi rời rạc; làm mờ một vùng có thể tạo ảnh không thuộc
phân phối dữ liệu mà mô hình từng gặp. Vì LIME học từ chính các phản ứng ấy,
phân phối phép nhiễu là một phần của lời giải thích, không phải chi tiết cài
đặt có thể bỏ qua.

